\section*{Hoofdstuk \ref{ch:g4:breathing} --- Ademhalen}

\begin{solution}{exo:g4:breathing:1}
\omterm{def:g4:breathing:movements}{Inademen}: de borst zet uit en er stroomt lucht naar binnen. \omterm{def:g4:breathing:movements}{Uitademen}:
de borst ontspant en de lucht stroomt naar buiten.
\end{solution}

\begin{solution}{exo:g4:breathing:2}
\omterm{def:g1:the-five-senses:sense}{Neus} (of \omterm{def:g4:journey-of-food:tract}{mond}), \omterm{def:g4:breathing:organs}{luchtpijp}, de twee takken, de \omterm{def:g4:breathing:organs}{longen}, en ten slotte de
miljoenen piepkleine luchtzakjes omwikkeld met bloedvaten.
\end{solution}

\begin{solution}{exo:g4:breathing:3}
Ongeveer $15$ keer per minuut. Bij een sprint stijgt dat naar $40$ of
meer, en ook dieper --- de hardwerkende \omterm{def:g3:bones-and-muscles:muscle}{spieren} hebben sneller zuurstof
nodig, dus laden de \omterm{def:g4:breathing:organs}{longen} die sneller.
\end{solution}

\begin{solution}{exo:g4:breathing:4}
Het neemt zuurstof op en geeft koolstofdioxide terug. De wisseling
gebeurt in de piepkleine luchtzakjes van de \omterm{def:g4:breathing:organs}{longen}, dwars door hun dunne
wanden, met het \omterm{prop:g4:journey-of-food:absorb}{bloed}.
\end{solution}

\begin{solution}{exo:g4:breathing:5}
Uitgeademde lucht is warmer, vochtiger (de beslagen spiegel), armer aan
zuurstof en rijker aan koolstofdioxide.
\end{solution}

\begin{solution}{exo:g4:breathing:6}
Omdat de wisseling tussen lucht en \omterm{prop:g4:journey-of-food:absorb}{bloed} oppervlak vraagt: miljoenen
zakjeswanden tellen op tot een halve tennisbaan aan ontmoetingsoppervlak,
opgevouwen in de borst. De \omterm{def:g4:journey-of-food:tract}{dunne darm} haalt met haar plooien dezelfde
truc uit om \omterm{def:g4:journey-of-food:digestion}{voedingsstoffen} op te nemen.
\end{solution}

\begin{solution}{exo:g4:breathing:7}
De \omterm{def:g1:the-five-senses:sense}{neus} verwarmt, bevochtigt en filtert de lucht onderweg naar binnen
--- koude, droge, stoffige lucht bereikt de tere \omterm{def:g4:breathing:organs}{longen} al voorbereid.
\end{solution}

\begin{solution}{exo:g4:breathing:8}
Open antwoord. Meestal ongeveer $15$ in rust en $25$--$40$ na de jumping
jacks: de werkende \omterm{def:g3:bones-and-muscles:muscle}{spieren} eisten meer zuurstof, dus versnelde en
verdiepte de ademhaling om die te leveren.
\end{solution}

\begin{solution}{exo:g4:breathing:9}
Dat het lichaam voedsel royaal opslaat (voor uren of dagen) maar bijna
helemaal geen voorraad zuurstof aanhoudt --- amper een minuut. Zuurstof
moet daarom voortdurend geleverd worden, teug na teug, en daarom kan de
ademhaling nooit lang pauzeren.
\end{solution}

\begin{solution}{exo:g4:breathing:10}
Vijfentwintig ademhalers hebben de hele ochtend zuurstof uit de lucht
van het lokaal gehaald en er koolstofdioxide voor teruggegeven: die
lucht is armer aan het ene en rijker aan het andere geworden. Het
voorschrift is regel 1: zet de ramen open en lucht het lokaal.
\end{solution}

\begin{solution}{exo:g4:breathing:11}
Allebei zijn ze de ontmoetingsplaats van \omterm{prop:g4:journey-of-food:absorb}{bloed} en zuurstof: enorme,
dunne, met \omterm{prop:g4:journey-of-food:absorb}{bloed} doorregen oppervlakken waar zuurstof naar binnen
oversteekt en koolstofdioxide naar buiten --- kieuwen die zuurstof
opgelost in water opnemen, \omterm{def:g4:breathing:organs}{longen} die haar uit de lucht halen. In de
lucht zakken de geveerde kieuwoppervlakken van de forel in elkaar en
plakken ze aan elkaar: het grote oversteekoppervlak krimpt tot bijna
niets, en zonder oversteekoppervlak geen zuurstof, hoe rijk de lucht
eromheen ook is.
\end{solution}
